% ===========================================================================
\chapter{Consolidated Numerical Results}
\label{app:results}
\lotgroup{Appendix A -- Consolidated Numerical Results}

\Cref{tab:experiments} lists the twelve principal experiments in chronological order.

\begin{table}[H]
\centering
\caption[The twelve principal experiments, in chronological order]{The twelve principal experiments, in chronological order.}
\label{tab:experiments}
\small
\begin{tabular}{@{}c P{4.6cm} l P{4.3cm} r@{}}
\toprule
\# & Experiment & Publication & Primary result & $N$ \\
\midrule
1 & GraphSAGE microservice orchestration & P1~\cite{belhadj2025icaart} & Cascade containment $81\%$ & 48 nodes \\
2 & Invariant-augmented GNN anomaly detection & P2, B1~\cite{belhadj2026invariants,belhadj2026lnaiinv} & F1 $94.4$--$97.0\%$ (P2); no significant F1 gain from invariants over ten runs (B1) & 4 datasets \\
3 & Hybrid CNN--GNN $+$ geometric and topological invariants & P3, B2~\cite{belhadj2026hybrid,mezghich2026lnaitopo} & $95.1\%$ binary accuracy ($+7.0$~pp) & 3{,}662 \\
4 & LightGBM-HC (hand-crafted, Messidor-2) & P4~\cite{belhadj2026icpram} & QWK $=0.858$, $1.6$~ms CPU classifier & 1{,}748 \\
5 & SAMF-GB & J2~\cite{belhadj2026samfgb} & QWK $=0.882$, $2.1$~ms CPU & 88{,}702 \\
6 & TOPORET on Kaggle DR & J3~\cite{belhadj2025plos} & QWK $=0.88$, accuracy $95.5\%$ & 88{,}702 \\
7 & TOPORET on Messidor-2 and APTOS 2019 & J3~\cite{belhadj2025plos} & QWK $=0.90$ and $0.86$ & 1{,}748; 3{,}662 \\
8 & H1 graph-contribution test & J3~\cite{belhadj2025plos} & $t_4=3.82$, $p=0.019$, $d_z=1.71$ & 88{,}702 \\
9 & H2 topological ANOVA & J3~\cite{belhadj2025plos} & $F=84.3$, $p<0.001$ & Kaggle DR \\
10 & DOTS in distribution & J1~\cite{belhadj2026dots} & QWK $=0.951$ & 88{,}702 \\
11 & DOTS H3a seven-criterion evaluation & J1~\cite{belhadj2026dots} & $7/7$ criteria met at the point estimate & 10{,}000 \\
12 & H3b necessity enumeration (investigated family) & J1~\cite{belhadj2026dots} & $6/6$ proper subsets fail & 10{,}000 \\
\bottomrule
\end{tabular}
\tabnote{$N$ is a number of images (or of graph nodes for experiments 1--2), not
necessarily of independent patients.}
\end{table}

\paragraph{Reading the sequence.}
The experiments are ordered chronologically rather than thematically, and the progression
is informative in itself. Experiments 1--2 evaluate the graph template entirely outside
medicine; the ten-run re-evaluation of experiment~2 revised its accuracy claim, leaving
interpretability as the main benefit of the invariants. Experiments 3--5 move to diabetic retinopathy at the proof-of-concept level,
each isolating one candidate signal against a fixed baseline. Experiments 6--9 form the
TOPORET evaluation, ending with the two confirmatory tests (8--9) that decide H1 and H2.
Experiments 10--12 apply the same discipline to DOTS and H3; experiment 12 is the most
computationally demanding of the programme, with six independently trained subset
architectures evaluated against all seven criteria under one protocol. The necessity enumeration (12) uses the $10{,}000$-image test set of the H3a evaluation, so
that DOTS and every proper subset are compared on identical data.


% ===========================================================================
\chapter{Hyperparameter Reference}
\label{app:hyper}
\lotgroup{Appendix B -- Hyperparameter Reference}

\Cref{tab:hyper} lists every hyperparameter of TOPORET and DOTS side by side.

\begin{table}[H]
\centering
\caption[Hyperparameter reference]{Hyperparameter reference: TOPORET versus DOTS.}
\label{tab:hyper}
\small
\begin{tabular}{@{}P{4.4cm}P{5.0cm}P{4.6cm}@{}}
\toprule
Hyperparameter & TOPORET~\cite{belhadj2025plos} & DOTS \\
\midrule
Backbone & EfficientNet-B3 (frozen) & RETFound ViT-L/16 (frozen) \\
Topological vector & 6-D (\cref{eq:tda-toporet}) & 7-D (\cref{eq:tda}) \\
Output head & Softmax, class-weighted CE & CORAL (architectural) \\
Graph construction & Threshold $\tau=0.65$, exact pairs & $k$-NN, $k=15$ (FAISS) \\
$\alpha$ (visual/topological balance) & 0.6 & 0.65 \\
Graph layer & GraphSAGE (mean) & Dynamic GAT \\
$d_h$ (hidden dimension) & 256 & 256 \\
$L$ (GNN layers) & 2 & 2 \\
$S$ (samples per layer) & 10 & 25 \\
Dropout $p$ & 0.30 & 0.30 \\
Class weights & $w_c=N/(C\,N_c)$ & $1.0, 2.5, 3.0, 8.0, 6.0$ \\
Label smoothing $\epsilon$ & -- & 0.05 \\
Abstention threshold $\tau_U$ & -- & 0.73~nats \\
Severe-DR referral threshold $t$ & -- & \refT{} (validation; frozen before test) \\
Referral rule & Not applicable & Refer if $\hat P(Y\geq3)>t$ or $H_i>\tau_U$ (\cref{eq:referral}) \\
Optimiser & Adam, lr $10^{-3}\to10^{-5}$ (cosine) & Adam, lr $=10^{-3}$ \\
Batch size & 64 & 32 \\
Epochs & $\leq200$, early stopping (20) & 100 \\
Data splits & 5 patient-aware $70/10/20$ & 5 folds, stratified \\
\bottomrule
\end{tabular}
\tabnote{The ordinal variant of TOPORET used as the starting point of DOTS
(\cref{ch:dots}) keeps the TOPORET architecture (dual-similarity GraphSAGE and topological
features) with a CORN head, and is trained under the DOTS protocol and graph settings.}
\end{table}

\paragraph{Reading the comparison.}
TOPORET and DOTS were developed under different protocols, and their hyperparameters
differ beyond the components that address Gaps~A--C (\cref{sec:ch4-gaps}): the
topological vector, the graph construction ($\tau$-threshold with exact pairwise
distances versus FAISS $k$-NN), the sampling depth, the optimisation schedule and the data
splits. The TOPORET-to-DOTS ablation of \cref{tab:ablate} is therefore performed within
the DOTS protocol, starting from the ordinal variant of TOPORET, so that each gap
resolution is introduced as an isolated change. The TOPORET hyperparameters were selected
on validation data (\cref{tab:hparam}).

% ===========================================================================
\chapter{Statistical Test Reference}
\label{app:stats}
\lotgroup{Appendix C -- Statistical Test Reference}

This appendix collects the formal statement of every statistical test used in the thesis.
Each hypothesis was pre-specified: before any experimental result was seen, the null
hypothesis, the test statistic, the significance threshold and the effect-size metric were
fixed. \Cref{tab:falsif} lists, for each claim, the outcome that would have falsified it,
the observed result and the margin by which the claim survived.

\begin{table}[H]
\centering
\caption[Falsification criteria and observed margins]{Falsification criteria and observed margins.}
\label{tab:falsif}
\small
\begin{tabular}{@{}P{4.2cm}P{3.0cm}P{4.3cm}P{2.9cm}@{}}
\toprule
Claim & Falsified if & Observed & Margin \\
\midrule
H1 -- graph QWK gain & $|t_4| < 2.78$ & $t_4=3.82$, $p=0.019$ & $\times1.37$ critical \\
H2 -- topology linked to severity & $F < 3.72$ & $F=84.3$, $p<0.001$ & $\times23$ critical \\
H3a -- all seven criteria & any criterion violated & $7/7$ met at the point estimate (\cref{tab:h3a}) & Sev-NR CI $[1.4\%,2.4\%]$ crosses $2\%$ \\
H3b -- necessity (investigated family) & any proper subset has Sev-NR $<2\%$ & $6/6$ proper subsets fail & constructive, scoped \\
CORAL violations & any violation at test time & $0.0\%$ & algebraic identity \\
Sev-NR point estimate & $\geq 2.0\%$ & $1.9\%$ $[1.4\%,\,2.4\%]$ & $0.1$~pp \\
\bottomrule
\end{tabular}
\tabnote{The critical values are $t_{0.025,4}=2.78$ (the conservative threshold adopted
for H1) and $F_{0.005;\,4,\,\infty}\approx3.72$ (conservative level used for H2); see
\cref{app:stats}.}
\end{table}

\paragraph{Notation convention.}
\textbf{H1}, \textbf{H2}, \textbf{H3} (upright) denote the three research hypotheses.
$H_0$ and $H_1$ (italic, subscript, in math mode) denote the zeroth and first
persistent-homology groups. The two notations are unrelated: hypothesis H2 is tested
using $H_0$ and $H_1$ features as independent variables.


\section{H1: Paired \texorpdfstring{$t$}{t}-Test}
\begin{align}
H_0 &: \mu_{\Delta\mathrm{QWK}} = 0 \quad \text{(the graph gives no QWK change)}, \\
H_a &: \mu_{\Delta\mathrm{QWK}} \neq 0 \quad \text{(the graph changes the QWK)}, \\
t &= \frac{\bar d}{s_d/\sqrt{n}}, \qquad \mathrm{df}=n-1=4,
\end{align}
where $\bar d$ is the mean of the five per-split QWK differences (TOPORET minus the CNN
baseline, over five patient-aware $70/10/20$ splits) and $s_d$ their standard deviation. Because the inferential unit is the
data split, the test has four degrees of freedom; the bootstrap and exact
permutation analyses of \cref{sec:testimpl} are complementary robustness checks on the
same five differences, not independent replications. The test is two-tailed, with
critical value $t_{0.025,4}=2.776$ at $\alpha=0.05$. Observed: $t_4=3.82$, two-tailed
$p=0.019$ (one-tailed $0.009$), on Kaggle DR; $t_4=4.1$ ($p=0.015$) on Messidor-2 and
$t_4=3.5$ ($p=0.025$) on APTOS~2019~\cite{belhadj2025plos}.

\section{H2: One-Way ANOVA}
\begin{align}
H_0 &: \mu_{G_0}=\mu_{G_1}=\cdots=\mu_{G_4} \quad \text{for total $H_1$ persistence } \Pi_{H_1}, \\
H_a &: \text{at least one } \mu_{G_k} \text{ differs}, \\
F &= \frac{\mathrm{MS}_{\mathrm{between}}}{\mathrm{MS}_{\mathrm{within}}}, \qquad
\mathrm{df}_1=4 .
\end{align}
A single primary feature, $\Pi_{H_1}$, is tested, at the conservative level $\alpha=0.005$;
for the large samples of Kaggle DR the critical value is $F_{0.005;\,4,\,\infty}\approx3.72$.
Post-hoc: Tukey HSD~\cite{tukey1949} on the pairwise grade comparisons. Observed:
$F=84.3$, $p<0.001$~\cite{belhadj2025plos}.

\section{H3a: Bootstrap Confidence Intervals}
For each of the seven criteria of \cref{tab:h3a}, a 95\% confidence interval is computed
by stratified bootstrap~\cite{efron1993bootstrap} ($B=10{,}000$ resamples, stratified by
true grade to preserve the class distribution). A criterion ``passes'' if its point
estimate meets the threshold; the interval is reported for transparency, but the
pre-specified decision rule uses the point estimate, except where explicitly stated (the
Phase~2 success criterion requires the interval).

\section{H3b: Constructive Enumeration}
No inferential test applies to H3b. The claim is established by constructing and
evaluating all six proper subsets against the same seven-criterion specification
(\cref{sec:h3b}). Claims of the form ``no member of this finite, fully enumerable set
satisfies property $P$'' are best established by exhaustive verification, since the whole
set is evaluated. The set, however, is the six subsets \emph{as implemented} in the
investigated family; the result does not extend to every conceivable architecture lacking
one of the axes, and each subset's Sev-NR is itself an estimate with sampling uncertainty. The price of this rigour is computational: every proper subset must be trained and evaluated on all benchmarks (\cref{tab:compute}).

\section{Effect-Size Formulae}
\begin{align}
d &= \frac{\bar x_1-\bar x_2}{s_{\mathrm{pooled}}}, \qquad
s_{\mathrm{pooled}} = \sqrt{\frac{(n_1-1)s_1^2+(n_2-1)s_2^2}{n_1+n_2-2}}, \\
d_z &= \frac{\bar\delta}{s_\delta} = \frac{t}{\sqrt{n}}, \\
\eta^2 &= \frac{\mathrm{SS}_{\mathrm{between}}}{\mathrm{SS}_{\mathrm{total}}}
        = \frac{df_1\,F}{df_1\,F+df_2} .
\end{align}
The first formula is used for comparisons between independent groups; the second, where
$\bar\delta$ and $s_\delta$ are the mean and standard deviation of the paired differences
over the $n$ splits, is used for the paired $t$-test of H1 ($d_z=3.82/\sqrt5=1.71$); the third
relates $\eta^2$ to the $F$ statistic of a one-way ANOVA, and requires the within-groups
degrees of freedom $df_2$, which are not reported with the H2 statistic
in~\cite{belhadj2025plos}.
Effect sizes complement $p$-values: a $p$-value says whether an effect is likely to be
real, an effect size says how large it is, independently of the sample size. With the large
samples used in this thesis, even negligible differences can reach significance, which is
why every primary comparison is reported with its effect size. \Cref{tab:effthresh} gives
the conventional interpretation thresholds~\cite{cohen1988}; all standardised effect sizes
reported in \cref{tab:effects} meet or exceed the ``large'' threshold.

\begin{table}[H]
\centering
\caption{Conventional thresholds for the interpretation of effect sizes.}
\label{tab:effthresh}
\small
\begin{tabular}{@{}lccc@{}}
\toprule
Measure & Small & Medium & Large \\
\midrule
Cohen's $d$ (difference of two means) & 0.2 & 0.5 & 0.8 \\
$\eta^2$ (share of variance explained, ANOVA) & 0.01 & 0.06 & 0.14 \\
\bottomrule
\end{tabular}
\end{table}

\section{Control of Multiple Comparisons}

Testing many hypotheses on the same data increases the chance of a false positive. Three
measures limit this risk in the thesis. First, each research hypothesis has a single
pre-specified primary test and a single primary metric, fixed before the experiment
(\cref{tab:falsif}); secondary analyses, such as the bootstrap and permutation tests of H1,
are reported as supporting evidence and do not change the decision. Second, when several
tests address the same hypothesis, as for the topological features of H2, the significance
level is corrected with the Bonferroni method and lowered further for conservatism. Third,
the necessity claim H3b does not rely on a statistical test at all, since the six proper
subsets form the complete set of alternatives within the investigated family and are all
evaluated.

% ===========================================================================
\chapter{Reproducibility and Software Environment}
\label{app:repro}
\lotgroup{Appendix D -- Reproducibility and Software Environment}

\section{Code and Data Availability}
All code, model weights and configuration files described in this thesis are released
under the Apache~2.0 licence at \url{https://github.com/Nader-BelHadj/plosene} upon
acceptance of the DOTS article~\cite{belhadj2026dots}. The release includes the source
code of the six pipeline stages (\cref{tab:stages}), the 47-test verification suite
(\cref{tab:unittests}), YAML configurations for the twelve principal experiments
(\cref{app:results}), a Dockerfile that pins every software dependency, and the
evaluation manifest \texttt{test\_manifest.csv} listing the identifier, source dataset and
grade of each image of the H3a test set (\cref{tab:testcomp}). The five
datasets are public and must be obtained from their original distributors
(\cref{app:datasets}).

\section{Software Environment}
\Cref{tab:software} lists the software components and their versions.
\begin{table}[H]
\centering
\caption[Software environment and key library versions]{Software environment and key library versions.}
\label{tab:software}
\small
\begin{tabular}{@{}l P{10.4cm}@{}}
\toprule
Component & Version / use \\
\midrule
Python & 3.10 \\
PyTorch~\cite{paszke2019pytorch} & 2.1, CUDA 12.1 \\
PyTorch Geometric~\cite{fey2019pyg} & 2.4 (GraphSAGE and GAT implementations) \\
GUDHI~\cite{maria2014gudhi} & 3.9 (Vietoris--Rips persistent homology, stage S3) \\
FAISS~\cite{johnson2021faiss} & 1.7 (approximate $k$-NN, stage S5) \\
timm & Version pinned in the Dockerfile (RETFound and EfficientNet-B3 loading) \\
scikit-learn~\cite{pedregosa2011sklearn} & 1.3 (evaluation metrics) \\
SciPy & ANOVA and $t$-tests (\cref{sec:testimpl}) \\
OpenCV & 4.8 (CLAHE and Zhang--Suen skeletonisation, stage S1) \\
\bottomrule
\end{tabular}
\end{table}

\Cref{tab:hardware} lists the hardware used at each stage of the programme. All timings reported in the thesis were measured on this hardware, so that latency and cost figures can be compared directly across chapters.

\begin{table}[!tb]
\centering
\caption[Hardware used for training, benchmarking and testing]{Hardware used for training, benchmarking and testing.}
\label{tab:hardware}
\small
\begin{tabular}{@{}l P{9.6cm}@{}}
\toprule
Context & Hardware \\
\midrule
Model training (all experiments) & $1\times$ NVIDIA A100 (40~GB), 16-core CPU, 128~GB RAM \\
Inference latency (\cref{tab:latency}) & $1\times$ NVIDIA A100 (stages S2, S5, S6); 16-core CPU (stage S3) \\
Unit-test suite & CPU only, any modern x86-64 workstation \\
\bottomrule
\end{tabular}
\end{table}

\section{Randomness and Seeding}
The primary experiments use a single global seed (42) for all stochastic operations: model
initialisation, data-loader shuffling, cross-validation fold assignment and bootstrap
resampling. The stability analysis of \cref{sec:numrepro} repeats the main experiments with two
additional seeds (123 and 7).
Five-fold splits are stratified by ICDR grade, computed once per experiment and held fixed
across all ablation variants (\cref{tab:ablate,tab:ablation3}), so that differences
between configurations are not confounded by different splits. The unit of every split and
of every reported sample size is the retinal image. Some benchmarks contain several images
of the same patient (for example both eyes in Kaggle EyePACS), so the reported results are
image-level estimates and patient-level generalisation is interpreted cautiously.

\section{Reproducibility Tolerance}
\label{sec:tolerance}
Because some GPU kernels are not bit-exact across hardware generations even with fixed
seeds, results within $2\sigma$ of the reported values -- with $\sigma$ the five-fold
standard deviation reported with each primary metric (for example \cref{tab:h1}) -- are
considered a successful reproduction. Bit-level reproduction is neither claimed nor
expected.

\section{FAIR Compliance Statement}
\begin{description}[style=nextline,leftmargin=1.2em]
\item[Findable] Code repository indexed with a persistent identifier on acceptance of the
DOTS article; thesis archived in the institutional repository of ENSI, University of
Manouba.
\item[Accessible] Public repository without access restrictions; the Apache~2.0 licence
permits reuse with attribution.
\item[Interoperable] Configurations in YAML; model weights in the standard PyTorch
checkpoint format.
\item[Reusable] Documented dependency versions (\cref{tab:software}), a Dockerfile for
exact environment reconstruction, and this appendix as a stand-alone reference.
\end{description}

\section{Computational Cost Summary}
\Cref{tab:compute} summarises approximate compute consumption by phase, from logged job
durations on the hardware of \cref{tab:hardware}. Figures are rounded, because several
exploratory runs of the early phases were not systematically logged. The H3b enumeration
is the most expensive phase despite testing the smallest architectures: exhaustive
enumeration requires every proper subset to be trained to convergence on every
benchmark, rather than sharing a backbone as the DOTS ablation (\cref{tab:ablate}) does.

\begin{table}[H]
\centering
\caption[Approximate GPU-hours by experimental phase]{Approximate GPU-hours by experimental phase.}
\label{tab:compute}
\small
\begin{tabular}{@{}P{6.2cm} r P{6.0cm}@{}}
\toprule
Phase & GPU-hours & Dominant cost \\
\midrule
Distributed-systems benchmarks (\cref{ch:found}) & $\sim$40 & SimPy simulation; CPU-bound \\
Structural proof of concept (\cref{ch:found}) & $\sim$120 & EfficientNet-B0 hybrid ablations; hand-crafted pipelines \\
TOPORET development and validation (\cref{ch:toporet}) & $\sim$310 & Pairwise Wasserstein graphs; hyperparameter sweep (\cref{tab:hparam}) $\times$ five splits \\
DOTS development and ablation (\cref{ch:dots}) & $\sim$480 & RETFound-based models across Gaps A--C, five folds \\
H3b necessity enumeration (\cref{ch:dots}) & $\sim$540 & Six subset architectures $\times$ five benchmarks \\
\midrule
Total & $\sim$1{,}490 & \\
\bottomrule
\end{tabular}
\end{table}

\section{Numerical Reproducibility Specification}
\label{sec:numrepro}
\Cref{tab:seeds} gives the seeds and the variance across runs for the main systems.

\begin{table}[H]
\centering
\caption[Reproducibility specification]{Reproducibility specification: mean $\pm$ SD over three seeds.}
\label{tab:seeds}
\small
\begin{tabular}{@{}llcc@{}}
\toprule
Experiment & Seeds & QWK & Sev-NR \\
\midrule
TOPORET, ordinal variant (Kaggle EyePACS) & 42, 123, 7 & $0.880\pm0.003$ & $3.4\pm0.2\%$ \\
DOTS (Kaggle EyePACS) & 42, 123, 7 & $0.951\pm0.003$ & $1.9\pm0.1\%$ \\
DOTS (APTOS 2019) & 42, 123, 7 & $0.953\pm0.004$ & $1.7\pm0.2\%$ \\
DOTS (IDRiD) & 42, 123, 7 & $0.878\pm0.005$ & $2.1\pm0.3\%$ \\
DOTS (Messidor-2) & 42, 123, 7 & $0.948\pm0.003$ & $1.8\pm0.2\%$ \\
DOTS (DDR) & 42, 123, 7 & $0.951\pm0.003$ & $1.9\pm0.2\%$ \\
\bottomrule
\end{tabular}
\end{table}

\section{Parameter Selection Record}
All parameters were selected on training and validation data only, in this order:
\begin{enumerate}
\item \textbf{Backbone} (EfficientNet-B3 or RETFound): validation-fold QWK; fixed before
the H3a test.
\item \textbf{TOPORET graph ($\alpha=0.6$, $\tau=0.65$)}: validation grid search
(\cref{tab:hparam}); fixed before the H1 test.
\item \textbf{DOTS graph ($\alpha=0.65$, $k=15$)}: validation-fold QWK; fixed before the H3a
test.
\item \textbf{Class weights $\lambda_0,\ldots,\lambda_4$}: grid search on validation-fold Sev-NR and macro-F1
(\cref{tab:weights}); fixed before the H3a test.
\item \textbf{Referral threshold $t=\refT$}: selected on the validation folds together with
$\tau_U$; fixed before the H3a test.
\item \textbf{Abstention threshold $\tau_U=0.73$ nats}: smallest threshold meeting Sev-NR
$<2\%$ with coverage above $90\%$ on the validation fold (\cref{tab:tau}); fixed before
the H3a test.
\item \textbf{Test-set access}: each held-out test set was accessed exactly once, after
all parameters above were fixed.
\end{enumerate}

\section{Statistical Test Implementation Details}
\label{sec:testimpl}
\begin{itemize}
\item \textbf{H1 $t$-test:} paired, two-tailed, df $=4$; threshold $|t_4|>2.776$; passed
($t_4=3.82$).
\item \textbf{H1 bootstrap:} $10{,}000$ resamples with replacement of the five split
differences; 95\% CI from the 2.5th and 97.5th percentiles.
\item \textbf{H1 permutation:} exact enumeration of the $2^5=32$ sign assignments of the
five split differences; the one-sided $p$-value is the fraction of assignments whose mean
difference is at least the observed one. With all five differences positive, $p=1/32$.
\item \textbf{H2 ANOVA:} one-way ANOVA of $\Pi_{H_1}$ across the five grades on Kaggle DR,
level $0.005$, Tukey HSD post-hoc~\cite{belhadj2025plos}.
\item \textbf{H3b bootstrap:} $10{,}000$ resamples of prediction-level Sev-NR for each
subset; the CI of the difference is computed as above.
\end{itemize}
